% ============================================================
% PREPARACION CONTROL -- CALCULO I
% Contenidos: paginas 128 a 135 de Apuntes Calculo I MBI
% ============================================================

\documentclass[letterpaper,12pt]{article}

\usepackage[utf8]{inputenc}
\usepackage[T1]{fontenc}
\usepackage[spanish]{babel}
\usepackage{amsmath,amssymb,amsfonts}
\usepackage{geometry}
\usepackage{fancyhdr}
\usepackage{enumitem}
\usepackage{graphicx}
\usepackage{hyperref}
\usepackage{multicol}
\usepackage{parskip}
\usepackage{xcolor}
\usepackage{tcolorbox}
\usepackage{eso-pic}
\makeatletter
\providecommand{\IfPDFManagementActiveF}[1]{#1}
\makeatother
\usepackage{transparent}
\usepackage{tikz}
\usepackage{needspace}
\tcbuselibrary{skins,breakable}

\geometry{letterpaper,margin=1.5cm,top=3cm,bottom=2.5cm,headheight=2cm}

\definecolor{celesteSuave}{HTML}{F0F7FF}
\definecolor{azulFuerte}{HTML}{1A5276}
\definecolor{azulMedio}{HTML}{2980B9}
\definecolor{verdeValidado}{HTML}{1E8449}
\definecolor{enlace}{HTML}{1A5276}

\newcommand{\logoup}{../../../../../template/images/logo_up.png}
\newcommand{\logousach}{../../../../../template/images/logo_usach.png}
\newcommand{\logofondo}{../../../../../template/images/logo_up.png}

\newcommand{\institucion}{Usach Premium}
\newcommand{\curso}{Cálculo I}
\newcommand{\guiasnumero}{Preparación Control}
\newcommand{\semestre}{1er. Semestre de 2026}
\newcommand{\preparadopor}{Jaime Riquelme}
\newcommand{\webinstitucion}{www.usachpremium.cl}
\newcommand{\instagraminstitucion}{https://www.instagram.com/usach.premium}
\newcommand{\transparencia}{0.045}
\newcommand{\fraseportada}{``Por y para estudiantes''}
\newcommand{\sen}{\operatorname{sen}}
\newcommand{\R}{\mathbb{R}}

\hypersetup{
    colorlinks=true,
    linkcolor=enlace,
    urlcolor=enlace,
    pdfborder={0 0 0},
    pdftitle={Preparación Control - Cálculo I},
    pdfauthor={Usach Premium}
}

\pagestyle{fancy}
\fancyhf{}
\fancyhead[L]{\includegraphics[height=1.2cm]{\logoup}}
\fancyhead[R]{%
    \begin{minipage}[b]{0.47\textwidth}
        \raggedleft
        \fontsize{9pt}{11pt}\selectfont
        \textbf{\color{azulFuerte}\institucion}\\
        \curso\ -- \guiasnumero\\
        \textit{\color{gray}@usach.premium}
    \end{minipage}\hspace{0.1cm}%
    \includegraphics[height=1.2cm]{\logousach}%
}
\fancyfoot[L]{\fontsize{9pt}{11pt}\selectfont\textit{\fraseportada}}
\fancyfoot[R]{\fontsize{9pt}{11pt}\selectfont Página \thepage}
\renewcommand{\headrulewidth}{0.4pt}
\renewcommand{\footrulewidth}{0.4pt}

\fancypagestyle{plain}{%
    \fancyhf{}
    \renewcommand{\headrulewidth}{0pt}
    \renewcommand{\footrulewidth}{0pt}
}

\newcommand{\MarcaAgua}{%
    \AddToShipoutPicture{%
        \AtPageCenter{%
            \makebox(0,0){%
                \transparent{\transparencia}%
                \includegraphics[width=0.7\textwidth]{\logofondo}%
            }%
        }%
    }%
}

\newtcolorbox{resumenbox}[1]{%
    colback=celesteSuave,
    colframe=azulFuerte,
    colbacktitle=azulFuerte,
    coltitle=white,
    title=\textbf{#1},
    fonttitle=\bfseries,
    arc=7pt,
    boxrule=1pt,
    enhanced,
    breakable
}

\newtcolorbox{controlbox}[1]{%
    colback=white,
    colframe=azulFuerte,
    colbacktitle=azulFuerte,
    coltitle=white,
    title=\textbf{#1},
    fonttitle=\bfseries,
    arc=7pt,
    boxrule=1pt,
    enhanced,
    drop shadow={black!6},
    breakable
}

\newtcolorbox{respuestasbox}[1]{%
    colback=celesteSuave,
    colframe=azulFuerte,
    title=\textbf{#1},
    fonttitle=\bfseries,
    arc=5pt,
    boxrule=0.9pt,
    breakable
}

\newtcolorbox{validacionbox}{%
    colback=verdeValidado!6,
    colframe=verdeValidado,
    coltext=black,
    arc=6pt,
    boxrule=1pt
}

\newcommand{\tituloSeccion}[1]{%
    \needspace{4\baselineskip}
    \subsection*{\color{azulFuerte}#1}
}

\setlist[enumerate]{itemsep=0.5em,topsep=0.5em}

\begin{document}
\hypersetup{pageanchor=false}
\thispagestyle{plain}

\AddToShipoutPicture*{%
    \put(54,700){\includegraphics[width=2cm]{\logoup}}
    \put(420,706){\includegraphics[width=5cm]{\logousach}}
}

\vspace*{0.5cm}

\begin{center}
    \color{azulFuerte}
    {\huge\textbf{GUÍA DE EJERCICIOS}}\\[0.5cm]
    {\Huge\textbf{\curso}}\\[0.3cm]
    {\Large\textbf{\guiasnumero}}\\[1.5cm]

    \begin{tcolorbox}[
        colback=celesteSuave,
        colframe=azulFuerte,
        arc=10pt,
        width=0.85\textwidth,
        center,
        boxrule=1.5pt]
        \centering
        \vspace{0.2cm}
        {\Large\textbf{Contenidos}}
        \vspace{0.3cm}
        \color{black}
        \begin{itemize}[leftmargin=2.1cm,itemsep=0.28cm]
            \item[\color{azulFuerte}$\Rightarrow$] Recta tangente y derivada por definición
            \item[\color{azulFuerte}$\Rightarrow$] Diferenciabilidad y límites laterales
            \item[\color{azulFuerte}$\Rightarrow$] Reglas de derivación
            \item[\color{azulFuerte}$\Rightarrow$] Regla de la cadena
            \item[\color{azulFuerte}$\Rightarrow$] 40 ejercicios con dificultad progresiva
        \end{itemize}
        \vspace{0.2cm}
    \end{tcolorbox}
\end{center}

\vspace{1.1cm}

\begin{center}
    {\Large\textbf{\semestre}}\\[0.7cm]
    {\large\textbf{\preparadopor}}
\end{center}

\vfill

\begin{center}
    {\color{azulFuerte}\hrule height 0.5pt}
    \vspace{0.3cm}
    \begin{minipage}{0.45\textwidth}
        \centering
        \textbf{Instagram:}\\
        \small\href{\instagraminstitucion}{@usach.premium}
    \end{minipage}
    \begin{minipage}{0.45\textwidth}
        \centering
        \textbf{Web:}\\
        \small\href{http://\webinstitucion}{\webinstitucion}
    \end{minipage}
\end{center}

\clearpage
\pagenumbering{arabic}
\hypersetup{pageanchor=true}
\MarcaAgua

\section*{\color{azulFuerte}Antes de comenzar}

\begin{resumenbox}{FÓRMULAS ESENCIALES}
\begin{multicols}{2}
\textbf{Definición y recta tangente}
\[
f'(x)=\lim_{h\to0}\frac{f(x+h)-f(x)}{h}.
\]
\[
y-f(x_0)=f'(x_0)(x-x_0).
\]

\columnbreak

\textbf{Reglas principales}
\[
(fg)'=f'g+fg',
\]
\[
\left(\frac{f}{g}\right)'=\frac{f'g-fg'}{g^2},
\]
\[
(f\circ g)'(x)=f'(g(x))g'(x).
\]
\end{multicols}
\end{resumenbox}

\begin{validacionbox}
\textbf{Sugerencia:} trabaja sin mirar las respuestas. Escribe la regla utilizada y conserva la forma factorizada cuando permita reconocer mejor la estructura de la función.
\end{validacionbox}

\section*{\color{azulFuerte}Ejercicios}

\tituloSeccion{I. Reglas básicas}
\noindent Halle la derivada de cada función en su dominio natural.

\begin{enumerate}[label=\color{azulFuerte}\textbf{\arabic*.},leftmargin=*]
    \item \(f(x)=5x^4-3x^2+7x-9.\)

    \item \(g(x)=\sqrt{x}+2\sen(x)-3\cos(x)+e^x+\ln(x).\)

    \item \(h(x)=(x^2+1)(3x-2).\)

    \item \(p(x)=\dfrac{2x+1}{x^2+4}.\)

    \item \(q(x)=x^3-\dfrac{\pi}{3}-x^2\cos(x).\)
\end{enumerate}

\tituloSeccion{II. Definición, diferenciabilidad y recta tangente}

\begin{controlbox}{6. DERIVADA POR DEFINICIÓN}
Use exclusivamente la definición de derivada para determinar \(f'(x)\), donde
\[
f(x)=x^2-4x+1.
\]
\end{controlbox}

\begin{controlbox}{7. DERIVADA POR DEFINICIÓN}
Considere la función
\[
f(x)=\frac{x}{x+2},
\]
definida en su dominio natural. Determine \(f'(x)\) mediante la definición.
\end{controlbox}

\begin{controlbox}{8. DIFERENCIABILIDAD}
Demuestre, mediante límites laterales del cociente incremental, que
\[
f(x)=|x-3|
\]
no es diferenciable en \(x=3\).
\end{controlbox}

\begin{controlbox}{9. RECTA TANGENTE}
Encuentre la ecuación de la recta tangente al gráfico de
\[
f(x)=-x^2+2x+2
\]
en el punto \((3,f(3))\).
\end{controlbox}

\begin{controlbox}{10. RECTA TANGENTE HORIZONTAL}
Considere
\[
f(x)=x+\frac{1}{x},
\]
definida en su dominio natural. Halle la ecuación de la recta tangente al gráfico de \(f\) en \(x=1\) e indique qué característica tiene dicha recta.
\end{controlbox}

\tituloSeccion{III. Producto, cociente y regla de la cadena}
\noindent Calcule la derivada de las siguientes funciones y simplifique cuando sea conveniente.

\begin{enumerate}[label=\color{azulFuerte}\textbf{\arabic*.},leftmargin=*,start=11,itemsep=1em]
    \item \(f(x)=(x^2-1)e^x.\)

    \item \(f(x)=\dfrac{\sen(x)}{x^2+1}.\)

    \item \(f(x)=(3-2x)^5.\)

    \item \(f(x)=\sen\!\left(\sqrt{2x+1}\right).\)

    \item \(f(x)=e^{x^2-3x}.\)

    \item \(f(x)=\ln\!\left((x^2+1)^3\right).\)
\end{enumerate}

\tituloSeccion{IV. Ejercicios tipo control}

\begin{controlbox}{17. OPERACIONES CON FUNCIONES DERIVABLES}
Sean \(f\) y \(g\) funciones derivables tales que
\[
f(2)=1,\qquad f'(2)=-2,\qquad g(2)=3,\qquad g'(2)=4.
\]
Determine:
\begin{enumerate}[label=\alph*),leftmargin=*]
    \item \((5f-2g)'(2)\).
    \item \((fg)'(2)\).
    \item \(\left(\dfrac{f}{g}\right)'(2)\).
\end{enumerate}
\end{controlbox}

\begin{controlbox}{18. PARÁMETRO Y RECTA TANGENTE}
Sea \(a\in\R\) y considere
\[
f_a(x)=\sqrt{ax+4}.
\]
Determine el valor de \(a\) para que la recta tangente al gráfico de \(f_a\) en \(x=0\) tenga pendiente \(3/4\). Luego, halle la ecuación de dicha recta.
\end{controlbox}

\begin{controlbox}{19. DIFERENCIABILIDAD CON PARÁMETROS}
Sean \(a,b\in\R\) y sea
\[
f(x)=
\begin{cases}
x^2+a, & x\leq1,\\
bx+1, & x>1.
\end{cases}
\]
Determine los valores de \(a\) y \(b\) para que \(f\) sea diferenciable en \(x=1\). Argumente las condiciones utilizadas.
\end{controlbox}

\begin{controlbox}{20. EJERCICIO ADAPTADO DE CONTROL ANTERIOR}
Considere la función
\[
f(x)=\frac{x^2-3}{x-2},
\]
definida en su dominio natural.
\begin{enumerate}[label=\alph*),leftmargin=*]
    \item Halle \(f'(x)\) y entregue el resultado simplificado.
    \item Encuentre la ecuación de la recta tangente al gráfico de \(f\) en \(x=3\).
\end{enumerate}
\end{controlbox}

\tituloSeccion{V. Segundo bloque: reglas de derivación}
\noindent Halle la derivada de cada función en su dominio natural.

\begin{enumerate}[label=\color{azulFuerte}\textbf{\arabic*.},leftmargin=*,start=21,itemsep=1em]
    \item \(f(x)=x^5-\dfrac{2}{x}+3\sqrt{x}.\)

    \item \(g(x)=(2x^2-3x+1)\sen(x).\)

    \item \(h(x)=\dfrac{e^x+1}{x-1}.\)
\end{enumerate}

\tituloSeccion{VI. Segundo bloque: definición y rectas tangentes}

\begin{controlbox}{24. DERIVADA POR DEFINICIÓN}
Considere la función \(f(x)=1/x\), definida en su dominio natural. Use exclusivamente la definición para determinar \(f'(x)\).
\end{controlbox}

\begin{controlbox}{25. DERIVADA POR DEFINICIÓN}
Sea \(f(x)=\sqrt{x}\). Mediante la definición y una racionalización adecuada, determine \(f'(x)\) para \(x>0\).
\end{controlbox}

\begin{controlbox}{26. DIFERENCIABILIDAD}
Demuestre, mediante límites laterales del cociente incremental, que
\[
f(x)=|2x+1|
\]
no es diferenciable en \(x=-1/2\).
\end{controlbox}

\begin{controlbox}{27. RECTA TANGENTE A UNA FUNCIÓN COMPUESTA}
Encuentre la ecuación de la recta tangente al gráfico de
\[
f(x)=\sqrt{4x+5}
\]
en el punto \((1,f(1))\).
\end{controlbox}

\begin{controlbox}{28. PENDIENTE PRESCRITA}
Considere la función
\[
f(x)=x^2-2x+3.
\]
Determine el punto de su gráfico donde la recta tangente tiene pendiente igual a \(4\) y halle la ecuación de dicha recta.
\end{controlbox}

\tituloSeccion{VII. Segundo bloque: regla de la cadena}
\noindent Calcule y simplifique la derivada de las siguientes funciones.

\begin{enumerate}[label=\color{azulFuerte}\textbf{\arabic*.},leftmargin=*,start=29,itemsep=1em]
    \item \(f(x)=(x^2-x+1)^{-7}.\)

    \item \(f(x)=\cos^3(2x-1).\)

    \item \(f(x)=\sqrt{3x^2+x+2}.\)

    \item \(f(x)=\ln\!\left(\dfrac{2x-1}{x+3}\right).\)

    \item \(f(x)=e^{\sqrt{x+1}}.\)

    \item \(f(x)=\sen(x)e^{x^2+1}.\)
\end{enumerate}

\tituloSeccion{VIII. Segundo bloque: ejercicios tipo PEP y control}

\begin{controlbox}{35. COMPOSICIÓN DE FUNCIONES DERIVABLES}
Sean \(f\) y \(g\) funciones derivables tales que
\[
f'(3)=-4,\qquad g(-1)=3,\qquad g'(-1)=5.
\]
Determine \((f\circ g)'(-1)\).
\end{controlbox}

\begin{controlbox}{36. PARÁMETRO Y PENDIENTE}
Sea \(a\in\R\) y considere
\[
f_a(x)=(ax-1)^4.
\]
Determine el valor de \(a\) para que la recta tangente al gráfico de \(f_a\) en \(x=0\) tenga pendiente \(8\). Luego, encuentre la ecuación de dicha recta.
\end{controlbox}

\begin{controlbox}{37. FUNCIÓN POR TRAMOS}
Sean \(a,b\in\R\) y sea
\[
f(x)=
\begin{cases}
x^2+ax+b, & x\leq0,\\
e^x, & x>0.
\end{cases}
\]
Determine los valores de \(a\) y \(b\) para que \(f\) sea diferenciable en \(x=0\). Argumente las condiciones utilizadas.
\end{controlbox}

\begin{controlbox}{38. RECTAS TANGENTES CON PENDIENTE DADA}
Considere
\[
f(x)=\frac{x}{x+1},
\]
definida en su dominio natural.
\begin{enumerate}[label=\alph*),leftmargin=*]
    \item Determine todos los puntos del gráfico donde la recta tangente tiene pendiente \(1/4\).
    \item Halle las ecuaciones de las rectas tangentes correspondientes.
\end{enumerate}
\end{controlbox}

\begin{controlbox}{39. REGLAS COMBINADAS}
Considere la función
\[
f(x)=\left(\frac{x^2+1}{x-1}\right)^3,
\]
definida en su dominio natural. Halle \(f'(x)\) y entregue el resultado simplificado.
\end{controlbox}

\begin{controlbox}{40. PROBLEMA INTEGRADO}
Considere la función
\[
f(x)=(2x-1)^4e^x.
\]
\begin{enumerate}[label=\alph*),leftmargin=*]
    \item Halle \(f'(x)\) y factorice el resultado.
    \item Encuentre la ecuación de la recta tangente al gráfico de \(f\) en \(x=0\).
\end{enumerate}
\end{controlbox}

\clearpage
\begin{center}
    \begin{tcolorbox}[
        colback=azulFuerte,
        colframe=azulFuerte,
        colupper=white,
        width=0.8\textwidth,
        center,
        arc=10pt]
        \centering\Large\textbf{RESPUESTAS FINALES}
    \end{tcolorbox}
\end{center}

\begin{respuestasbox}{I. Reglas básicas}
\begin{enumerate}[label=\color{azulFuerte}\textbf{\arabic*.},leftmargin=*,itemsep=0.8em]
    \item \(\boxed{f'(x)=20x^3-6x+7}\).
    \item \(\boxed{g'(x)=\dfrac{1}{2\sqrt{x}}+2\cos(x)+3\sen(x)+e^x+\dfrac1x}\), \(x>0\).
    \item \(\boxed{h'(x)=9x^2-4x+3}\).
    \item \(\boxed{p'(x)=\dfrac{-2x^2-2x+8}{(x^2+4)^2}}\).
    \item \(\boxed{q'(x)=3x^2-2x\cos(x)+x^2\sen(x)}\).
\end{enumerate}
\end{respuestasbox}

\begin{respuestasbox}{II. Definición, diferenciabilidad y recta tangente}
\begin{enumerate}[label=\color{azulFuerte}\textbf{\arabic*.},leftmargin=*,start=6,itemsep=0.8em]
    \item \(\boxed{f'(x)=2x-4}\).
    \item \(\boxed{f'(x)=\dfrac{2}{(x+2)^2}}\), \(x\neq-2\).
    \item Los límites laterales del cociente incremental son \(\boxed{-1}\) y \(\boxed{1}\); por tanto, \(f\) no es diferenciable en \(x=3\).
    \item \(\boxed{y=-4x+11}\).
    \item \(\boxed{y=2}\); la recta tangente es horizontal.
\end{enumerate}
\end{respuestasbox}

\begin{respuestasbox}{III. Producto, cociente y regla de la cadena}
\begin{enumerate}[label=\color{azulFuerte}\textbf{\arabic*.},leftmargin=*,start=11,itemsep=0.8em]
    \item \(\boxed{f'(x)=e^x(x^2+2x-1)}\).
    \item \(\boxed{f'(x)=\dfrac{(x^2+1)\cos(x)-2x\sen(x)}{(x^2+1)^2}}\).
    \item \(\boxed{f'(x)=-10(3-2x)^4}\).
    \item \(\boxed{f'(x)=\dfrac{\cos(\sqrt{2x+1})}{\sqrt{2x+1}}}\), \(x>-\tfrac12\).
    \item \(\boxed{f'(x)=(2x-3)e^{x^2-3x}}\).
    \item \(\boxed{f'(x)=\dfrac{6x}{x^2+1}}\).
\end{enumerate}
\end{respuestasbox}

\begin{respuestasbox}{IV. Ejercicios tipo control}
\begin{enumerate}[label=\color{azulFuerte}\textbf{\arabic*.},leftmargin=*,start=17,itemsep=0.9em]
    \item \(\boxed{\text{a) }-18\qquad \text{b) }-2\qquad \text{c) }-\dfrac{10}{9}}\).
    \item \(\boxed{a=3}\) y \(\boxed{y=\dfrac34x+2}\).
    \item \(\boxed{a=2,\quad b=2}\).
    \item \(\boxed{\text{a) }f'(x)=\dfrac{x^2-4x+3}{(x-2)^2},\ x\neq2\qquad \text{b) }y=6}\).
\end{enumerate}
\end{respuestasbox}

\begin{respuestasbox}{V. Segundo bloque: reglas de derivación}
\begin{enumerate}[label=\color{azulFuerte}\textbf{\arabic*.},leftmargin=*,start=21,itemsep=0.8em]
    \item \(\boxed{f'(x)=5x^4+\dfrac{2}{x^2}+\dfrac{3}{2\sqrt{x}}}\), \(x>0\).
    \item \(\boxed{g'(x)=(4x-3)\sen(x)+(2x^2-3x+1)\cos(x)}\).
    \item \(\boxed{h'(x)=\dfrac{e^x(x-2)-1}{(x-1)^2}}\), \(x\neq1\).
\end{enumerate}
\end{respuestasbox}

\begin{respuestasbox}{VI. Segundo bloque: definición y rectas tangentes}
\begin{enumerate}[label=\color{azulFuerte}\textbf{\arabic*.},leftmargin=*,start=24,itemsep=0.8em]
    \item \(\boxed{f'(x)=-\dfrac{1}{x^2}}\), \(x\neq0\).
    \item \(\boxed{f'(x)=\dfrac{1}{2\sqrt{x}}}\), \(x>0\).
    \item Los límites laterales son \(\boxed{-2}\) y \(\boxed{2}\); por tanto, \(f\) no es diferenciable en \(x=-1/2\).
    \item \(\boxed{y=\dfrac23x+\dfrac73}\).
    \item Punto: \(\boxed{(3,6)}\). Recta tangente: \(\boxed{y=4x-6}\).
\end{enumerate}
\end{respuestasbox}

\begin{respuestasbox}{VII. Segundo bloque: regla de la cadena}
\begin{enumerate}[label=\color{azulFuerte}\textbf{\arabic*.},leftmargin=*,start=29,itemsep=0.8em]
    \item \(\boxed{f'(x)=-7(2x-1)(x^2-x+1)^{-8}}\).
    \item \(\boxed{f'(x)=-6\cos^2(2x-1)\sen(2x-1)}\).
    \item \(\boxed{f'(x)=\dfrac{6x+1}{2\sqrt{3x^2+x+2}}}\).
    \item \(\boxed{f'(x)=\dfrac{7}{(2x-1)(x+3)}}\), \(x\in(-\infty,-3)\cup(1/2,+\infty)\).
    \item \(\boxed{f'(x)=\dfrac{e^{\sqrt{x+1}}}{2\sqrt{x+1}}}\), \(x>-1\).
    \item \(\boxed{f'(x)=e^{x^2+1}\bigl(\cos(x)+2x\sen(x)\bigr)}\).
\end{enumerate}
\end{respuestasbox}

\begin{respuestasbox}{VIII. Segundo bloque: ejercicios tipo PEP y control}
\begin{enumerate}[label=\color{azulFuerte}\textbf{\arabic*.},leftmargin=*,start=35,itemsep=0.9em]
    \item \(\boxed{(f\circ g)'(-1)=-20}\).
    \item \(\boxed{a=-2}\) y \(\boxed{y=8x+1}\).
    \item \(\boxed{a=1,\quad b=1}\).
    \item Puntos: \(\boxed{(1,1/2)}\) y \(\boxed{(-3,3/2)}\). Rectas: \(\boxed{y=\dfrac14x+\dfrac14}\) y \(\boxed{y=\dfrac14x+\dfrac94}\).
    \item \(\boxed{f'(x)=\dfrac{3(x^2+1)^2(x^2-2x-1)}{(x-1)^4}}\), \(x\neq1\).
    \item \(\boxed{\text{a) }f'(x)=e^x(2x-1)^3(2x+7)\qquad \text{b) }y=-7x+1}\).
\end{enumerate}
\end{respuestasbox}

\vspace{0.6cm}
\begin{center}
    \small\textbf{Usach Premium} -- Nos vemos en la ayudantía.\\
    \textit{``Por y para estudiantes.''}
\end{center}

\end{document}
